\documentclass[12pt]{article}
\usepackage{amsmath, amssymb, amsthm}
\usepackage{geometry}
\usepackage{physics}
\usepackage{bm}
\usepackage{graphicx}
\usepackage{hyperref}

\geometry{margin=1in}

\title{AFET: Entropically Coupled Threshold Dynamics \\
Formal Extension of UTAC as an Information–System Interaction Framework}

\author{}
\date{}

\begin{document}

\maketitle

\begin{abstract}
We formalize the Adaptive Field–Entropic Threshold (AFET) framework as an entropically coupled extension of the Universal Threshold Activation Criticality (UTAC) model. 
UTAC provides a nonlinear information-theoretic activation structure. AFET introduces entropy production as a necessary dynamical constraint governing parameter evolution within UTAC. 
We demonstrate that entropy production acts as a modulation operator over threshold position, coupling steepness, and damping. 
The resulting formulation describes a hierarchical, domain-dependent threshold system capable of modeling multi-scale dynamical transitions.
\end{abstract}

\section{Foundations}

\subsection{UTAC: Nonlinear Threshold Activation}

UTAC defines the activation probability of a system state transition as:

\begin{equation}
A(t) = \sigma\left( \beta (R - \Theta) \right)
\end{equation}

with logistic activation:

\begin{equation}
\sigma(x) = \frac{1}{1 + e^{-x}}
\end{equation}

Parameters:

\begin{itemize}
\item $R$ — system response variable
\item $\Theta$ — activation threshold
\item $\beta$ — coupling/steepness parameter
\item $A$ — activation state probability
\end{itemize}

The width of the transition region is:

\begin{equation}
\Delta R \propto \frac{1}{\beta}
\end{equation}

Thus $\beta$ controls sensitivity and critical sharpness.

\section{Entropy as Dynamical Driver}

We define thermodynamic entropy:

\begin{equation}
S = k_B \ln \Omega
\end{equation}

Information-theoretic entropy:

\begin{equation}
H = - \sum_i p_i \ln p_i
\end{equation}

Entropy production rate (crucial quantity):

\begin{equation}
\Pi = \frac{dS}{dt}
\end{equation}

Key principle:

\begin{quote}
No entropy production ($\Pi = 0$) implies no sustained dynamical state transitions.
\end{quote}

Thus entropy production acts as an enabling condition for UTAC dynamics.

\section{AFET: Entropically Coupled UTAC}

AFET introduces entropy-dependent modulation of UTAC parameters.

\subsection{General Coupled Form}

\begin{equation}
A(t) = \sigma\left(
\beta_{\text{eff}}(t) 
\left(
R(t) - \Theta_{\text{eff}}(t)
\right)
\right)
\end{equation}

with:

\begin{align}
\beta_{\text{eff}}(t) &= \beta_0 + f_\beta(\Pi, \kappa_{\text{int}}, \kappa_{\text{ext}}) \\
\Theta_{\text{eff}}(t) &= \Theta_0 + f_\Theta(\Pi, \rho) \\
\zeta_{\text{eff}}(t) &= \zeta_0 + f_\zeta(\Pi)
\end{align}

where:

\begin{itemize}
\item $\Pi$ — entropy production rate
\item $\kappa$ — internal/external coupling strength
\item $\rho$ — system order parameter
\item $\zeta$ — damping coefficient
\end{itemize}

\section{Hierarchical Beta Structure}

We define $\beta$ as a hierarchical stochastic field:

\begin{equation}
\beta_{d,t,i}(t) =
\mu_d
+ \delta_{d,t}
+ \delta_{d,t,i}
+ g(\Pi_{d,t,i}(t))
+ \epsilon
\end{equation}

Indices:

\begin{itemize}
\item $d$ — domain
\item $t$ — system type
\item $i$ — individual system
\end{itemize}

This defines:

\begin{itemize}
\item Domain-level band
\item System-type band
\item System-specific variation
\item State-dependent drift
\end{itemize}

\section{Necessary Conditions for Dynamical Universality}

A functioning dynamical system requires:

\begin{equation}
\Pi > 0
\end{equation}

and

\begin{equation}
\exists \; \beta, \Theta \text{ such that } \frac{dA}{dR} \neq 0
\end{equation}

Thus:

\begin{equation}
\text{Dynamical viability} \iff
(\Pi > 0) \land (\text{Nonlinear threshold structure})
\end{equation}

\section{Falsifiability Criteria}

The AFET framework is falsified if:

\begin{enumerate}
\item No measurable covariance exists between entropy production $\Pi$ and threshold dynamics.
\item $\beta$ remains invariant across domain perturbations.
\item Transition widths remain constant under dissipative shifts.
\item Systems sustain transitions at $\Pi \to 0$.
\end{enumerate}

\section{Conclusion}

UTAC defines informational activation structure.
AFET defines entropic governance over activation parameters.

Entropy production enables state change.
UTAC structures state change.

Together they form a minimal dynamical architecture capable of modeling multi-domain threshold systems.

\end{document}
